\documentclass[11pt]{article}
\usepackage{amsmath,amssymb}
\title{Room 25 (seat: direct): checking the ``$2\pi$ safety margin / net-loss when
$\ell\ge2$'' condition at $N=6$'s actual parameters}
\date{2026-09-27}
\begin{document}
\maketitle

\section*{0. What was read}

\texttt{N6\_full\_status\_v2.tex} (Reviewer BF's gap description); \texttt{room21\_seat\_direct.tex}
in full (domination constant derivation, \S3); \texttt{code/zeta\_probe/rootunity\_zeros/
qi\_bounds.tex}, the printed \texttt{lem:qi-bounds} statement and, crucially, the proof of part
(ii) around the sentence containing ``net loss'' (the line right after the boxed Poisson/triple-
product identity in the proof block, not the lemma statement itself -- the phrase does not appear
in the statement, only in the proof).

\section*{1. The exact statement, extracted}

The relevant sentence, verbatim from the proof of part (ii) in \texttt{qi\_bounds.tex}:

\begin{quote}
``The index maximising the leading part is one of the two class members next to the vertex
$m^*(x)$. Any other index loses $\ge c\,l(l-1)$ and gains at most $\frac\pi2|2-j|\,l$ in the
linear part. For $r\le\frac\pi2\cos\theta$ this is a net loss when $l\ge2$, which accounts for the
terms $\frac\pi2(|m^*|+4)$ and $2\pi$ in $\Lambda'$.''
\end{quote}

\textbf{What is being compared to what.} This is \emph{not} about the ``$2\pi$'' bookkeeping
identity ($N\times2\pi/N=2\pi$) that Reviewer BF correctly flagged as a tautology -- that identity
is a separate, trivial fact about the sum $\sum_j\frac\pi2|2-j|$ ($=\pi(1+0+1+2)=4\pi$, or the
related quantity). The sentence above is a genuine inequality with content: it bounds, for a
\emph{fixed} residue class $j$ (equivalently a fixed $m$-class mod $N$), the cost of using the
$l$-th-nearest class member (in the $k$-index, $l=|k-k_*|$) instead of the best one, when
assembling $\max_{m\equiv\varrho}$ into the sum that becomes $U(x)$ plus the $\Lambda'$/$K'$ error
terms.
\begin{itemize}
\item \textbf{Loss}: the quadratic part of $A_k$ (the exponent, \S3 of \texttt{room21}) is concave
in $k$ with leading coefficient $-c$, so stepping $l$ places from the optimal $k_*$ costs at least
$c\,l(l-1)$ in the exponent -- this is exactly the domination-lemma quadratic used to bound
$\sum_ke^{A_k}\le e^{A_{k_*}}(3+3e^{-2c})$.
\item \textbf{Gain}: the \emph{linear}-in-$y$ term of $\beta_k^2/(2Nt)$, coefficient
$\frac{2(N/2-j)}N$ (room21 \S2's boxed formula), moves by $\pi$ per unit $l$-step (since $m$ shifts
by $N$ per unit $l$ and $y=x-\frac{i\pi m}N$ shifts by $i\pi l$), so the linear part changes by at
most $\bigl|\frac{2(N/2-j)}N\bigr|\cdot\pi l$ as $l$ grows -- worst case over $j$ gives the stated
$\frac\pi2|2-j|l$ at $N=4$.
\item \textbf{Condition}: loss dominates gain once $c\,l(l-1)\ge(\text{gain coefficient})\cdot l$,
i.e.\ (worst case $l=2$, the first non-trivial step) $c\ge(\text{gain coefficient})$. This is what
``net loss when $\ell\ge2$'' means: for $l\ge2$ the quadratic loss already exceeds the linear gain,
so the two class-neighbors nearest $m^*$ dominate everything else, letting the sum over classes be
replaced by the finite $\Lambda'$ correction terms instead of a genuinely infinite worst case.
\end{itemize}

\section*{2. The general-$N$ form of the condition, derived}

\textbf{Gain coefficient, general $N$.} Using room21 \S2's $N$-generic linear coefficient
$\frac{2(N/2-j)}N$ and the $\pi$-per-unit-$l$ shift in $y$ (this shift rate is itself
$N$-independent: $m$ steps by $N$ per unit $l$, and $y=x-i\pi m/N$, so the shift in $y$ is
$i\pi N l/N=i\pi l$, cancelling the $N$), the gain per unit $l$ for class $j$ is
\[
  \frac{2\pi}N\Bigl|\frac N2-j\Bigr|\,l,\qquad j=1,\dots,N.
\]
The worst class is $j=N$ (or $j=1$), where $|N/2-j|=N/2$, giving gain $\le\frac{2\pi}N\cdot\frac
N2\cdot l=\pi l$ -- \textbf{exactly $\pi l$ for every $N$}, matching $N=4$'s stated
$\frac\pi2|2-j|l\le\pi l$ (at $j=4$, $|2-4|=2$, $\frac\pi2\cdot2=\pi$) exactly. This is a genuine,
non-tautological $N$-independent fact: the worst-case linear gain per unit $l$-step is $\pi$
regardless of $N$, because the $(N/2-j)$ range and the $1/N$ prefactor cancel.

\textbf{Loss coefficient, general $N$.} From room21 \S3, $c_N=\dfrac{2\pi^2\cos\theta}{Nr}$ (this
is the $c$ in $A_{k_*+l}\le A_{k_*}-c_N\,l(l-1)$; checked $c_N|_{N=4}=\pi^2\cos\theta/(2r)$, the
printed value).

\textbf{The condition, general $N$.} Net loss at $l=2$ (the binding case, since $l(l-1)/l=l-1$ is
increasing in $l$, so $l=2$ is the tightest constraint) requires
\[
  c_N\ \ge\ \pi \quad\Longleftrightarrow\quad \frac{2\pi^2\cos\theta}{Nr}\ge\pi
  \quad\Longleftrightarrow\quad \boxed{r\ \le\ \frac{2\pi\cos\theta}N}.
\]
\textbf{Check at $N=4$}: $r\le\frac{2\pi\cos\theta}4=\frac\pi2\cos\theta$ -- \emph{exactly} the
printed condition ``$r\le\frac\pi2\cos\theta$''. This confirms the extraction and the
generalization method (same cross-check discipline as room21).

\textbf{At $N=6$}: $r\le\dfrac{2\pi\cos\theta}6=\dfrac{\pi\cos\theta}3$.

\section*{3. Plugging in $N=6$'s actual parameters}

At the certified saddle angle $\theta^*=0$ (Room 18/19), $\cos\theta=1$, so the threshold is
\[
  r_{\max}=\frac\pi3\approx1.0472.
\]
For the operating radius: two values appear in the room files with slightly different status.
\texttt{qi\_bounds.tex}'s own setting line fixes $\bar r=0.005$ as the \emph{global} radius cap
for the whole $N=4$ analysis (``$0<r\le\bar r=0.005$''); this is the value \texttt{N6\_full\_status\_v2.tex}
attributes to Room 21, but \textbf{Room 21 itself never writes $0.005$} -- its own \S6 numerical
check uses $r=0.003$ explicitly as ``the actual operating regime'' for $N=6$. This is a minor
mis-citation in the status file (conflating qi\_bounds's $N=4$ constant $\bar r=0.005$ with an
$N=6$-specific value), not a computational error, and it does not matter for the verdict: both
candidate values are tiny next to $\pi/3$.
\[
  r=0.005:\quad \frac{r_{\max}}r=\frac{\pi/3}{0.005}\approx209,
  \qquad
  r=0.003:\quad \frac{r_{\max}}r=\frac{\pi/3}{0.003}\approx349.
\]

\section*{4. Verdict}

\textbf{CONDITION HOLDS, by a very large margin (roughly 200--350$\times$), at $N=6$'s actual
operating parameters.} The condition $r\le\frac{2\pi\cos\theta}N$ (general-$N$ form of
\texttt{qi\_bounds.tex}'s ``$r\le\frac\pi2\cos\theta$, net loss when $\ell\ge2$'') at $N=6$,
$\theta=\theta^*=0$ becomes $r\le\pi/3\approx1.047$, and $N=6$'s certified/reported operating
radius ($r\approx0.003$--$0.005$) sits roughly $200$--$350$ times below that threshold. This is
not marginal in any sense -- the weaker domination exponent found in Room 21
($2\pi^2\cos\theta/(3r)$ vs.\ $N=4$'s $\pi^2\cos\theta/(2r)$, a factor $3/2$ smaller) does not
come close to threatening this particular inequality, because the threshold radius itself scales
down by the same $1/N$ factor as the domination constant's inverse-$r$ scaling, and $N=6$'s actual
$r$ is four orders of magnitude below either threshold.

\textbf{What this closes, precisely.} This resolves the \emph{net-loss-condition} half of
Reviewer BF's second gap: the concern (``could plausibly fail'') is checked and does not fail, by
a wide margin, at $N=6$'s certified saddle. It does \emph{not} by itself certify that the resulting
$\Lambda'$/$K'$ bookkeeping constants ($\frac\pi2(|m^*|+4)\to$ its $N=6$ analogue, already computed
in room21 \S7 as $\frac\pi2(|m^*|+9)$, and the ``$2\pi$'' term) are otherwise correctly assembled
into a working $N=6$ lemma -- that assembly still needs the missing $f_e^{(6)},f_o^{(6)}$ closed
forms identified as blocked in room21 \S7, and $c_\ell^{(6)}$ (Reviewer BF's \emph{first} gap,
untouched by this room). Those remain open.

\end{document}
